%%%%%%%%%%%%%%%%%%%%%%%%%%%%%%%%%%%%%%%%%%%%%%%%%%%%%%%%%%%%
% 5.17 ASYMPTOTIC ANALYSIS
% The Composition of Advanced Mathematics
%%%%%%%%%%%%%%%%%%%%%%%%%%%%%%%%%%%%%%%%%%%%%%%%%%%%%%%%%%%%

\section{Asymptotic Analysis}
\label{sec:asymptotic-analysis}

\prerequisite{
Calculus, real analysis, complex analysis, differential equations,
Fourier analysis, and approximation theory.
}

\learningobjectives{
By the end of this section, the reader should be able to:
\begin{itemize}
    \item explain the purpose of asymptotic analysis;
    \item use Big-O, little-o, and asymptotic equivalence notation;
    \item compare the relative growth of functions;
    \item distinguish convergent expansions from asymptotic expansions;
    \item manipulate asymptotic expansions algebraically;
    \item derive asymptotic approximations from Taylor expansions;
    \item understand asymptotic sequences and scales;
    \item approximate integrals using Laplace's method;
    \item apply Watson's lemma;
    \item understand the method of steepest descent;
    \item apply stationary-phase approximations;
    \item derive Stirling's approximation;
    \item understand dominant balance in differential equations;
    \item recognize regular and singular perturbation problems;
    \item construct boundary-layer approximations;
    \item understand matched asymptotic expansions;
    \item recognize multiple-scale methods;
    \item understand asymptotic inversion and implicit approximations;
    \item connect asymptotic analysis with special functions,
          probability, differential equations, and mathematical physics.
\end{itemize}
}

%%%%%%%%%%%%%%%%%%%%%%%%%%%%%%%%%%%%%%%%%%%%%%%%%%%%%%%%%%%%
\subsection{What Is Asymptotic Analysis?}
\label{subsec:what-is-asymptotic-analysis}
%%%%%%%%%%%%%%%%%%%%%%%%%%%%%%%%%%%%%%%%%%%%%%%%%%%%%%%%%%%%

Asymptotic analysis studies the behavior of mathematical quantities as a
parameter approaches a limiting regime.

Typical limits include

\[
x\to\infty,
\qquad
x\to0,
\qquad
n\to\infty,
\qquad
\varepsilon\to0^+.
\]

Rather than finding an exact formula, one seeks a simpler expression that
captures the dominant behavior.

For example,

\[
\sqrt{x^2+x}
\]

is complicated compared with \(x\), but as \(x\to\infty\),

\[
\sqrt{x^2+x}
=
x\sqrt{1+\frac1x}
\sim x.
\]

Thus

\[
\boxed{
\sqrt{x^2+x}\sim x
\qquad
(x\to\infty).
}
\]

The symbol \(\sim\), read ``is asymptotic to,'' expresses asymptotic
equivalence.

%%%%%%%%%%%%%%%%%%%%%%%%%%%%%%%%%%%%%%%%%%%%%%%%%%%%%%%%%%%%
\subsection{Asymptotic Equivalence}
\label{subsec:asymptotic-equivalence}
%%%%%%%%%%%%%%%%%%%%%%%%%%%%%%%%%%%%%%%%%%%%%%%%%%%%%%%%%%%%

\begin{definitionbox}{Asymptotic Equivalence}
Let \(f\) and \(g\) be functions with \(g(x)\neq0\) near the limiting
regime.

We write

\[
\boxed{
f(x)\sim g(x)
\qquad
(x\to a)
}
\]

if

\[
\boxed{
\lim_{x\to a}
\frac{f(x)}{g(x)}
=
1.
}
\]
\end{definitionbox}

Thus asymptotic equivalence means that the relative error tends to zero.

Indeed,

\[
\frac{f-g}{g}
=
\frac{f}{g}-1
\to0.
\]

%%%%%%%%%%%%%%%%%%%%%%%%%%%%%%%%%%%%%%%%%%%%%%%%%%%%%%%%%%%%
\subsection{Examples of Asymptotic Equivalence}
\label{subsec:asymptotic-equivalence-examples}
%%%%%%%%%%%%%%%%%%%%%%%%%%%%%%%%%%%%%%%%%%%%%%%%%%%%%%%%%%%%

As \(x\to0\),

\[
\boxed{
\sin x\sim x.
}
\]

Also,

\[
\boxed{
1-\cos x\sim\frac{x^2}{2}.
}
\]

As \(x\to\infty\),

\[
\boxed{
x^2+3x+1\sim x^2.
}
\]

As \(n\to\infty\),

\[
\boxed{
n+1\sim n.
}
\]

These relations isolate the dominant behavior.

%%%%%%%%%%%%%%%%%%%%%%%%%%%%%%%%%%%%%%%%%%%%%%%%%%%%%%%%%%%%
\subsection{Big-O Notation}
\label{subsec:asymptotic-big-o}
%%%%%%%%%%%%%%%%%%%%%%%%%%%%%%%%%%%%%%%%%%%%%%%%%%%%%%%%%%%%

\begin{definitionbox}{Big-O Notation}
We write

\[
\boxed{
f(x)=O(g(x))
\qquad
(x\to a)
}
\]

if there exist constants \(C>0\) and a neighborhood of the limiting
regime such that

\[
\boxed{
|f(x)|
\leq
C|g(x)|.
}
\]
\end{definitionbox}

Big-O describes an upper bound on magnitude.

For example,

\[
\boxed{
3x^2+5x+1
=
O(x^2)
\qquad
(x\to\infty).
}
\]

%%%%%%%%%%%%%%%%%%%%%%%%%%%%%%%%%%%%%%%%%%%%%%%%%%%%%%%%%%%%
\subsection{Little-o Notation}
\label{subsec:asymptotic-little-o}
%%%%%%%%%%%%%%%%%%%%%%%%%%%%%%%%%%%%%%%%%%%%%%%%%%%%%%%%%%%%

\begin{definitionbox}{Little-o Notation}
We write

\[
\boxed{
f(x)=o(g(x))
\qquad
(x\to a)
}
\]

if

\[
\boxed{
\frac{f(x)}{g(x)}
\to0.
}
\]
\end{definitionbox}

Thus \(f\) is asymptotically negligible compared with \(g\).

For example,

\[
\boxed{
x=o(x^2)
\qquad
(x\to\infty).
}
\]

Also,

\[
\boxed{
x^2=o(x)
\qquad
(x\to0).
}
\]

%%%%%%%%%%%%%%%%%%%%%%%%%%%%%%%%%%%%%%%%%%%%%%%%%%%%%%%%%%%%
\subsection{Big-O Versus Little-o}
\label{subsec:big-o-versus-little-o}
%%%%%%%%%%%%%%%%%%%%%%%%%%%%%%%%%%%%%%%%%%%%%%%%%%%%%%%%%%%%

The statement

\[
f=O(g)
\]

means that \(f\) is no larger than a constant multiple of \(g\) in the
relevant limiting regime.

The stronger statement

\[
f=o(g)
\]

means that \(f\) becomes negligible relative to \(g\).

Therefore,

\[
\boxed{
f=o(g)
\Longrightarrow
f=O(g),
}
\]

under the usual local assumptions.

The converse is false.

%%%%%%%%%%%%%%%%%%%%%%%%%%%%%%%%%%%%%%%%%%%%%%%%%%%%%%%%%%%%
\subsection{Asymptotic Equivalence and Little-o}
\label{subsec:equivalence-little-o}
%%%%%%%%%%%%%%%%%%%%%%%%%%%%%%%%%%%%%%%%%%%%%%%%%%%%%%%%%%%%

The relation

\[
f\sim g
\]

is equivalent to

\[
\boxed{
f=g+o(g).
}
\]

Indeed,

\[
f\sim g
\]

means

\[
\frac{f-g}{g}\to0.
\]

This is precisely

\[
f-g=o(g).
\]

%%%%%%%%%%%%%%%%%%%%%%%%%%%%%%%%%%%%%%%%%%%%%%%%%%%%%%%%%%%%
\subsection{Growth Hierarchies}
\label{subsec:asymptotic-growth-hierarchies}
%%%%%%%%%%%%%%%%%%%%%%%%%%%%%%%%%%%%%%%%%%%%%%%%%%%%%%%%%%%%

As \(x\to\infty\), a standard hierarchy is

\[
\boxed{
1
\ll
\log x
\ll
x^\alpha
\ll
a^x
\ll
x!
}
\]

for

\[
\alpha>0,
\qquad
a>1.
\]

Here

\[
f\ll g
\]

informally means that \(f\) grows much more slowly than \(g\).

More precisely, one often expresses this as

\[
\boxed{
f=o(g).
}
\]

For example,

\[
\boxed{
\log x=o(x^\alpha)
}
\]

for every \(\alpha>0\).

%%%%%%%%%%%%%%%%%%%%%%%%%%%%%%%%%%%%%%%%%%%%%%%%%%%%%%%%%%%%
\subsection{Algebra of Big-O Terms}
\label{subsec:big-o-algebra}
%%%%%%%%%%%%%%%%%%%%%%%%%%%%%%%%%%%%%%%%%%%%%%%%%%%%%%%%%%%%

Suppose

\[
f=O(g)
\]

and

\[
h=O(k).
\]

Then, under appropriate assumptions,

\[
\boxed{
f+h
=
O(|g|+|k|),
}
\]

and

\[
\boxed{
fh
=
O(gk).
}
\]

If

\[
f=O(x^m),
\qquad
h=O(x^n),
\]

then

\[
\boxed{
fh=O(x^{m+n}).
}
\]

These rules allow error terms to be manipulated systematically.

%%%%%%%%%%%%%%%%%%%%%%%%%%%%%%%%%%%%%%%%%%%%%%%%%%%%%%%%%%%%
\subsection{Asymptotic Expansion}
\label{subsec:asymptotic-expansion}
%%%%%%%%%%%%%%%%%%%%%%%%%%%%%%%%%%%%%%%%%%%%%%%%%%%%%%%%%%%%

An asymptotic expansion represents a function using a sequence of terms
of progressively smaller size.

Suppose

\[
\phi_0(x),
\phi_1(x),
\phi_2(x),\ldots
\]

satisfy

\[
\phi_{n+1}(x)=o(\phi_n(x)).
\]

Then one may write

\[
\boxed{
f(x)
\sim
\sum_{n=0}^{\infty}
a_n\phi_n(x)
}
\]

if, for every fixed \(N\),

\[
\boxed{
f(x)
-
\sum_{n=0}^{N}
a_n\phi_n(x)
=
o(\phi_N(x)).
}
\]

%%%%%%%%%%%%%%%%%%%%%%%%%%%%%%%%%%%%%%%%%%%%%%%%%%%%%%%%%%%%
\subsection{Asymptotic Sequence}
\label{subsec:asymptotic-sequence}
%%%%%%%%%%%%%%%%%%%%%%%%%%%%%%%%%%%%%%%%%%%%%%%%%%%%%%%%%%%%

\begin{definitionbox}{Asymptotic Sequence}
A sequence of functions

\[
\phi_0,\phi_1,\phi_2,\ldots
\]

is an asymptotic sequence in a given limiting regime if

\[
\boxed{
\phi_{n+1}=o(\phi_n)
}
\]

for every \(n\).
\end{definitionbox}

A common example as \(x\to\infty\) is

\[
\boxed{
1,
\frac1x,
\frac1{x^2},
\frac1{x^3},
\ldots
}
\]

because

\[
\frac{x^{-(n+1)}}{x^{-n}}
=
\frac1x
\to0.
\]

%%%%%%%%%%%%%%%%%%%%%%%%%%%%%%%%%%%%%%%%%%%%%%%%%%%%%%%%%%%%
\subsection{Convergent Series Versus Asymptotic Series}
\label{subsec:convergent-versus-asymptotic-series}
%%%%%%%%%%%%%%%%%%%%%%%%%%%%%%%%%%%%%%%%%%%%%%%%%%%%%%%%%%%%

An asymptotic expansion need not converge.

This is a fundamental distinction.

A convergent series seeks

\[
\boxed{
f
=
\sum_{n=0}^{\infty}a_n.
}
\]

An asymptotic series instead requires only that truncating after a finite
number of terms gives progressively better approximations in the
specified limiting regime.

Thus

\[
\boxed{
\text{asymptotic usefulness}
\neq
\text{infinite-series convergence}.
}
\]

%%%%%%%%%%%%%%%%%%%%%%%%%%%%%%%%%%%%%%%%%%%%%%%%%%%%%%%%%%%%
\subsection{Taylor Series as Asymptotic Expansions}
\label{subsec:taylor-asymptotic-expansions}
%%%%%%%%%%%%%%%%%%%%%%%%%%%%%%%%%%%%%%%%%%%%%%%%%%%%%%%%%%%%

Taylor expansions provide basic examples.

As \(x\to0\),

\[
e^x
=
1+x+\frac{x^2}{2}
+
O(x^3).
\]

Similarly,

\[
\sin x
=
x-\frac{x^3}{6}
+
O(x^5),
\]

and

\[
\cos x
=
1-\frac{x^2}{2}
+
O(x^4).
\]

Even when one does not require convergence of the full Taylor series,
finite Taylor expansions provide asymptotic information.

%%%%%%%%%%%%%%%%%%%%%%%%%%%%%%%%%%%%%%%%%%%%%%%%%%%%%%%%%%%%
\subsection{Worked Example: Elementary Expansion}
\label{subsec:worked-elementary-asymptotic}
%%%%%%%%%%%%%%%%%%%%%%%%%%%%%%%%%%%%%%%%%%%%%%%%%%%%%%%%%%%%

\begin{workedexamplebox}{Expansion of \(\sqrt{1+x}\)}

For \(x\to0\), use the binomial expansion:

\[
(1+x)^{1/2}
=
1+\frac12x
+
\frac{\frac12(-\frac12)}{2}x^2
+
O(x^3).
\]

Therefore,

\[
(1+x)^{1/2}
=
1+\frac{x}{2}
-\frac{x^2}{8}
+
O(x^3).
\]

\begin{answerbox}
\[
\boxed{
\sqrt{1+x}
=
1+\frac{x}{2}
-\frac{x^2}{8}
+
O(x^3).
}
\]
\end{answerbox}

\end{workedexamplebox}

%%%%%%%%%%%%%%%%%%%%%%%%%%%%%%%%%%%%%%%%%%%%%%%%%%%%%%%%%%%%
\subsection{Expansion of Quotients}
\label{subsec:asymptotic-quotients}
%%%%%%%%%%%%%%%%%%%%%%%%%%%%%%%%%%%%%%%%%%%%%%%%%%%%%%%%%%%%

Suppose

\[
f(x)
=
1+a_1x+a_2x^2+O(x^3)
\]

and

\[
g(x)
=
1+b_1x+b_2x^2+O(x^3).
\]

To expand

\[
\frac{f(x)}{g(x)},
\]

first use

\[
\frac1{1+u}
=
1-u+u^2+O(u^3).
\]

Then multiply the resulting expansions and retain terms only through the
required order.

Careful order bookkeeping is essential.

%%%%%%%%%%%%%%%%%%%%%%%%%%%%%%%%%%%%%%%%%%%%%%%%%%%%%%%%%%%%
\subsection{Asymptotic Inversion}
\label{subsec:asymptotic-inversion}
%%%%%%%%%%%%%%%%%%%%%%%%%%%%%%%%%%%%%%%%%%%%%%%%%%%%%%%%%%%%

Suppose

\[
y
=
x+a+\frac{b}{x}
+
O(x^{-2})
\]

as \(x\to\infty\).

One may seek an inverse expansion

\[
x
=
y+c+\frac{d}{y}
+
O(y^{-2}).
\]

Substituting this proposed expansion into the original relation and
matching orders determines the unknown coefficients.

This procedure is called asymptotic inversion.

%%%%%%%%%%%%%%%%%%%%%%%%%%%%%%%%%%%%%%%%%%%%%%%%%%%%%%%%%%%%
\subsection{Implicit Asymptotic Expansions}
\label{subsec:implicit-asymptotic-expansions}
%%%%%%%%%%%%%%%%%%%%%%%%%%%%%%%%%%%%%%%%%%%%%%%%%%%%%%%%%%%%

Many quantities are defined implicitly rather than by explicit formulas.

Suppose

\[
F(x,\varepsilon)=0
\]

and

\[
\varepsilon\to0.
\]

One may propose

\[
\boxed{
x(\varepsilon)
=
x_0+\varepsilon x_1+\varepsilon^2x_2+\cdots
}
\]

and substitute into the equation.

Matching equal powers of \(\varepsilon\) produces equations for

\[
x_0,x_1,x_2,\ldots.
\]

%%%%%%%%%%%%%%%%%%%%%%%%%%%%%%%%%%%%%%%%%%%%%%%%%%%%%%%%%%%%
\subsection{Dominant Balance}
\label{subsec:dominant-balance}
%%%%%%%%%%%%%%%%%%%%%%%%%%%%%%%%%%%%%%%%%%%%%%%%%%%%%%%%%%%%

Dominant balance identifies which terms control an equation in a limiting
regime.

Consider

\[
x^3+x=N
\]

for large \(N\).

If \(x\) is large, then

\[
x^3
\]

dominates \(x\).

Thus the leading balance is

\[
x^3\sim N.
\]

Hence

\[
\boxed{
x\sim N^{1/3}.
}
\]

This leading approximation can then be refined.

%%%%%%%%%%%%%%%%%%%%%%%%%%%%%%%%%%%%%%%%%%%%%%%%%%%%%%%%%%%%
\subsection{Scaling}
\label{subsec:asymptotic-scaling}
%%%%%%%%%%%%%%%%%%%%%%%%%%%%%%%%%%%%%%%%%%%%%%%%%%%%%%%%%%%%

Scaling introduces variables whose magnitude remains order one in a
limiting regime.

Suppose

\[
x\sim\varepsilon^\alpha.
\]

Define

\[
\boxed{
X=\frac{x}{\varepsilon^\alpha}.
}
\]

Then

\[
X=O(1).
\]

Choosing the correct scaling often reveals the dominant terms of a
problem.

%%%%%%%%%%%%%%%%%%%%%%%%%%%%%%%%%%%%%%%%%%%%%%%%%%%%%%%%%%%%
\subsection{Distinguished Limits}
\label{subsec:distinguished-limits}
%%%%%%%%%%%%%%%%%%%%%%%%%%%%%%%%%%%%%%%%%%%%%%%%%%%%%%%%%%%%

When several small or large parameters occur, different terms may
compete.

A distinguished limit is a scaling relationship between parameters that
keeps multiple important effects in balance.

For example, if

\[
\varepsilon x
\]

must remain order one as

\[
\varepsilon\to0,
\]

then the relevant scale is

\[
\boxed{
x=O(\varepsilon^{-1}).
}
\]

Distinguished limits are central in perturbation theory.

%%%%%%%%%%%%%%%%%%%%%%%%%%%%%%%%%%%%%%%%%%%%%%%%%%%%%%%%%%%%
\subsection{Asymptotic Approximation of Integrals}
\label{subsec:asymptotic-integrals}
%%%%%%%%%%%%%%%%%%%%%%%%%%%%%%%%%%%%%%%%%%%%%%%%%%%%%%%%%%%%

A major class of problems concerns integrals depending on a large
parameter.

Typical examples have the form

\[
\boxed{
I(\lambda)
=
\int_a^b
e^{\lambda\phi(x)}
\psi(x)\,dx
}
\]

with

\[
\lambda\to\infty.
\]

The exponential factor strongly favors regions where \(\phi\) is largest.

This observation leads to Laplace's method.

%%%%%%%%%%%%%%%%%%%%%%%%%%%%%%%%%%%%%%%%%%%%%%%%%%%%%%%%%%%%
\subsection{Laplace's Method}
\label{subsec:laplaces-method}
%%%%%%%%%%%%%%%%%%%%%%%%%%%%%%%%%%%%%%%%%%%%%%%%%%%%%%%%%%%%

Suppose \(\phi\) has a unique interior maximum at \(x_0\) with

\[
\phi'(x_0)=0
\]

and

\[
\phi''(x_0)<0.
\]

Near \(x_0\),

\[
\phi(x)
\approx
\phi(x_0)
+
\frac12\phi''(x_0)(x-x_0)^2.
\]

Then, under suitable hypotheses,

\[
\boxed{
\int_a^b
e^{\lambda\phi(x)}
\psi(x)\,dx
\sim
e^{\lambda\phi(x_0)}
\psi(x_0)
\sqrt{
\frac{2\pi}
{\lambda|\phi''(x_0)|}
}.
}
\]

This is the leading term of Laplace's method.

%%%%%%%%%%%%%%%%%%%%%%%%%%%%%%%%%%%%%%%%%%%%%%%%%%%%%%%%%%%%
\subsection{Why Laplace's Method Works}
\label{subsec:why-laplace-method-works}
%%%%%%%%%%%%%%%%%%%%%%%%%%%%%%%%%%%%%%%%%%%%%%%%%%%%%%%%%%%%

If

\[
\phi(x)<\phi(x_0)
\]

away from \(x_0\), then

\[
e^{\lambda\phi(x)}
\]

is exponentially smaller away from the maximum.

Therefore the integral becomes concentrated near \(x_0\).

The local quadratic approximation of \(\phi\) produces a Gaussian
integral, explaining the appearance of

\[
\boxed{
\sqrt{\frac{2\pi}{\lambda|\phi''(x_0)|}}.
}
\]

%%%%%%%%%%%%%%%%%%%%%%%%%%%%%%%%%%%%%%%%%%%%%%%%%%%%%%%%%%%%
\subsection{Worked Example: Laplace's Method}
\label{subsec:worked-laplace-method}
%%%%%%%%%%%%%%%%%%%%%%%%%%%%%%%%%%%%%%%%%%%%%%%%%%%%%%%%%%%%

\begin{workedexamplebox}{Gaussian-Type Integral}

Consider

\[
I(\lambda)
=
\int_{-1}^{1}
e^{-\lambda x^2}\,dx,
\qquad
\lambda\to\infty.
\]

Here

\[
\phi(x)=-x^2
\]

has its maximum at

\[
x_0=0.
\]

Also,

\[
\phi(0)=0,
\qquad
\phi''(0)=-2.
\]

Laplace's method gives

\[
I(\lambda)
\sim
\sqrt{
\frac{2\pi}
{2\lambda}
}.
\]

Therefore,

\begin{answerbox}
\[
\boxed{
I(\lambda)
\sim
\sqrt{\frac{\pi}{\lambda}}.
}
\]
\end{answerbox}

\end{workedexamplebox}

%%%%%%%%%%%%%%%%%%%%%%%%%%%%%%%%%%%%%%%%%%%%%%%%%%%%%%%%%%%%
\subsection{Endpoint Contributions}
\label{subsec:laplace-endpoint-contributions}
%%%%%%%%%%%%%%%%%%%%%%%%%%%%%%%%%%%%%%%%%%%%%%%%%%%%%%%%%%%%

The dominant contribution to an integral need not come from an interior
critical point.

If the maximum of \(\phi\) occurs at an endpoint, a different local
expansion is required.

For example,

\[
\int_0^b
e^{-\lambda x}\psi(x)\,dx
\]

is concentrated near

\[
x=0
\]

as

\[
\lambda\to\infty.
\]

This leads naturally to Watson's lemma.

%%%%%%%%%%%%%%%%%%%%%%%%%%%%%%%%%%%%%%%%%%%%%%%%%%%%%%%%%%%%
\subsection{Watson's Lemma}
\label{subsec:watsons-lemma}
%%%%%%%%%%%%%%%%%%%%%%%%%%%%%%%%%%%%%%%%%%%%%%%%%%%%%%%%%%%%

Suppose a function has a local expansion near \(t=0\),

\[
f(t)
\sim
\sum_{n=0}^{\infty}
a_nt^{n+\alpha-1},
\qquad
\alpha>0.
\]

Under suitable conditions,

\[
\int_0^\infty
e^{-xt}f(t)\,dt
\]

has an expansion as \(x\to\infty\):

\[
\boxed{
\int_0^\infty
e^{-xt}f(t)\,dt
\sim
\sum_{n=0}^{\infty}
a_n
\Gamma(n+\alpha)
x^{-(n+\alpha)}.
}
\]

Here \(\Gamma\), read ``Gamma,'' denotes the Gamma function.

Watson's lemma is especially useful for Laplace-type integrals.

%%%%%%%%%%%%%%%%%%%%%%%%%%%%%%%%%%%%%%%%%%%%%%%%%%%%%%%%%%%%
\subsection{The Gamma Function}
\label{subsec:asymptotic-gamma-function}
%%%%%%%%%%%%%%%%%%%%%%%%%%%%%%%%%%%%%%%%%%%%%%%%%%%%%%%%%%%%

The Gamma function is defined for suitable \(z\) by

\[
\boxed{
\Gamma(z)
=
\int_0^\infty
t^{z-1}e^{-t}\,dt.
}
\]

For positive integers,

\[
\boxed{
\Gamma(n+1)=n!.
}
\]

Thus asymptotic analysis of the Gamma function also gives asymptotic
information about factorials.

%%%%%%%%%%%%%%%%%%%%%%%%%%%%%%%%%%%%%%%%%%%%%%%%%%%%%%%%%%%%
\subsection{Stirling's Approximation}
\label{subsec:stirlings-approximation}
%%%%%%%%%%%%%%%%%%%%%%%%%%%%%%%%%%%%%%%%%%%%%%%%%%%%%%%%%%%%

One of the most important asymptotic formulas is

\[
\boxed{
n!
\sim
\sqrt{2\pi n}
\left(\frac{n}{e}\right)^n.
}
\]

Equivalently,

\[
\boxed{
\Gamma(x+1)
\sim
\sqrt{2\pi x}
\left(\frac{x}{e}\right)^x
}
\]

as \(x\to\infty\).

This is Stirling's approximation.

%%%%%%%%%%%%%%%%%%%%%%%%%%%%%%%%%%%%%%%%%%%%%%%%%%%%%%%%%%%%
\subsection{Logarithmic Form of Stirling's Formula}
\label{subsec:log-stirling}
%%%%%%%%%%%%%%%%%%%%%%%%%%%%%%%%%%%%%%%%%%%%%%%%%%%%%%%%%%%%

Taking logarithms gives

\[
\boxed{
\log(n!)
=
n\log n-n
+\frac12\log(2\pi n)
+
o(1).
}
\]

A refined expansion is

\[
\boxed{
\log(n!)
=
n\log n-n
+\frac12\log(2\pi n)
+\frac{1}{12n}
+O(n^{-3}).
}
\]

Logarithmic forms are often easier to use when factorials become very
large.

%%%%%%%%%%%%%%%%%%%%%%%%%%%%%%%%%%%%%%%%%%%%%%%%%%%%%%%%%%%%
\subsection{Deriving Stirling's Formula by Laplace's Method}
\label{subsec:stirling-laplace}
%%%%%%%%%%%%%%%%%%%%%%%%%%%%%%%%%%%%%%%%%%%%%%%%%%%%%%%%%%%%

Starting from

\[
\Gamma(n+1)
=
\int_0^\infty
t^ne^{-t}\,dt,
\]

set

\[
t=ns.
\]

Then

\[
\Gamma(n+1)
=
n^{n+1}
\int_0^\infty
e^{n(\log s-s)}\,ds.
\]

The phase

\[
\phi(s)=\log s-s
\]

has its maximum at

\[
s=1.
\]

Laplace's method then produces the leading Stirling approximation.

%%%%%%%%%%%%%%%%%%%%%%%%%%%%%%%%%%%%%%%%%%%%%%%%%%%%%%%%%%%%
\subsection{Oscillatory Integrals}
\label{subsec:asymptotic-oscillatory-integrals}
%%%%%%%%%%%%%%%%%%%%%%%%%%%%%%%%%%%%%%%%%%%%%%%%%%%%%%%%%%%%

Another major class of asymptotic problems has the form

\[
\boxed{
I(\lambda)
=
\int_a^b
e^{i\lambda\phi(x)}
a(x)\,dx.
}
\]

Unlike Laplace-type integrals, the magnitude of the exponential is

\[
|e^{i\lambda\phi(x)}|=1.
\]

Decay comes from cancellation caused by rapid oscillation.

%%%%%%%%%%%%%%%%%%%%%%%%%%%%%%%%%%%%%%%%%%%%%%%%%%%%%%%%%%%%
\subsection{Nonstationary Phase}
\label{subsec:nonstationary-phase}
%%%%%%%%%%%%%%%%%%%%%%%%%%%%%%%%%%%%%%%%%%%%%%%%%%%%%%%%%%%%

If

\[
\phi'(x)\neq0
\]

throughout the region of integration, integration by parts produces
decay in \(\lambda\).

Indeed,

\[
e^{i\lambda\phi(x)}
=
\frac{1}{i\lambda\phi'(x)}
\frac{d}{dx}
e^{i\lambda\phi(x)}.
\]

Repeated integration by parts can yield

\[
\boxed{
I(\lambda)=O(\lambda^{-N})
}
\]

for arbitrarily large \(N\), provided sufficient smoothness and boundary
conditions are available.

%%%%%%%%%%%%%%%%%%%%%%%%%%%%%%%%%%%%%%%%%%%%%%%%%%%%%%%%%%%%
\subsection{Stationary Phase}
\label{subsec:stationary-phase}
%%%%%%%%%%%%%%%%%%%%%%%%%%%%%%%%%%%%%%%%%%%%%%%%%%%%%%%%%%%%

The strongest contributions to oscillatory integrals often occur where

\[
\boxed{
\phi'(x_0)=0.
}
\]

Such points are called stationary points.

Suppose

\[
\phi''(x_0)\neq0.
\]

Then a leading stationary-phase approximation has the form

\[
\boxed{
I(\lambda)
\sim
a(x_0)
e^{i\lambda\phi(x_0)}
e^{i\frac{\pi}{4}\operatorname{sgn}(\phi''(x_0))}
\sqrt{
\frac{2\pi}
{\lambda|\phi''(x_0)|}
}.
}
\]

Here \(\operatorname{sgn}\) denotes the sign function.

%%%%%%%%%%%%%%%%%%%%%%%%%%%%%%%%%%%%%%%%%%%%%%%%%%%%%%%%%%%%
\subsection{Steepest Descent}
\label{subsec:steepest-descent}
%%%%%%%%%%%%%%%%%%%%%%%%%%%%%%%%%%%%%%%%%%%%%%%%%%%%%%%%%%%%

For complex contour integrals of the form

\[
\boxed{
I(\lambda)
=
\int_{\mathcal C}
e^{\lambda\phi(z)}
a(z)\,dz,
}
\]

one may deform the contour through saddle points of \(\phi\).

A saddle point satisfies

\[
\boxed{
\phi'(z_0)=0.
}
\]

The contour is chosen so that

\[
\operatorname{Re}\phi(z)
\]

decreases as rapidly as possible away from the saddle.

This is the method of steepest descent.

%%%%%%%%%%%%%%%%%%%%%%%%%%%%%%%%%%%%%%%%%%%%%%%%%%%%%%%%%%%%
\subsection{Saddle-Point Method}
\label{subsec:saddle-point-method}
%%%%%%%%%%%%%%%%%%%%%%%%%%%%%%%%%%%%%%%%%%%%%%%%%%%%%%%%%%%%

The saddle-point method is a broad complex-analytic framework encompassing
many Laplace and steepest-descent approximations.

The central steps are:

\begin{enumerate}
    \item locate critical points of the exponent;
    \item determine which critical points contribute;
    \item deform the contour when analytically permissible;
    \item expand the exponent locally;
    \item evaluate the resulting Gaussian-type integral;
    \item estimate the remaining error.
\end{enumerate}

It is widely used for special functions, combinatorics, probability, and
statistical mechanics.

%%%%%%%%%%%%%%%%%%%%%%%%%%%%%%%%%%%%%%%%%%%%%%%%%%%%%%%%%%%%
\subsection{Perturbation Theory}
\label{subsec:asymptotic-perturbation-theory}
%%%%%%%%%%%%%%%%%%%%%%%%%%%%%%%%%%%%%%%%%%%%%%%%%%%%%%%%%%%%

A perturbation problem contains a small parameter

\[
\varepsilon.
\]

One seeks an approximation to a solution

\[
u(x;\varepsilon)
\]

as

\[
\varepsilon\to0.
\]

A common assumption is a power expansion

\[
\boxed{
u
\sim
u_0
+
\varepsilon u_1
+
\varepsilon^2u_2
+\cdots.
}
\]

Substituting this expansion into the governing equation and matching
powers of \(\varepsilon\) produces a sequence of simpler problems.

%%%%%%%%%%%%%%%%%%%%%%%%%%%%%%%%%%%%%%%%%%%%%%%%%%%%%%%%%%%%
\subsection{Regular Perturbation}
\label{subsec:regular-perturbation}
%%%%%%%%%%%%%%%%%%%%%%%%%%%%%%%%%%%%%%%%%%%%%%%%%%%%%%%%%%%%

A perturbation is regular when the limiting problem obtained by setting

\[
\varepsilon=0
\]

retains the essential structure of the original problem.

For example,

\[
u'(x)+u(x)=\varepsilon f(x)
\]

has the same differential order when

\[
\varepsilon=0.
\]

A regular expansion

\[
u=u_0+\varepsilon u_1+\cdots
\]

is often valid throughout the domain.

%%%%%%%%%%%%%%%%%%%%%%%%%%%%%%%%%%%%%%%%%%%%%%%%%%%%%%%%%%%%
\subsection{Worked Example: Regular Perturbation}
\label{subsec:worked-regular-perturbation}
%%%%%%%%%%%%%%%%%%%%%%%%%%%%%%%%%%%%%%%%%%%%%%%%%%%%%%%%%%%%

\begin{workedexamplebox}{Algebraic Perturbation}

Consider

\[
x+\varepsilon x^2=1,
\qquad
\varepsilon\to0.
\]

Assume

\[
x
=
x_0+\varepsilon x_1+O(\varepsilon^2).
\]

Substitute:

\[
x_0+\varepsilon x_1
+
\varepsilon x_0^2
+
O(\varepsilon^2)
=
1.
\]

At order \(1\),

\[
x_0=1.
\]

At order \(\varepsilon\),

\[
x_1+x_0^2=0.
\]

Hence

\[
x_1=-1.
\]

Therefore,

\begin{answerbox}
\[
\boxed{
x
=
1-\varepsilon
+
O(\varepsilon^2).
}
\]
\end{answerbox}

\end{workedexamplebox}

%%%%%%%%%%%%%%%%%%%%%%%%%%%%%%%%%%%%%%%%%%%%%%%%%%%%%%%%%%%%
\subsection{Singular Perturbation}
\label{subsec:singular-perturbation}
%%%%%%%%%%%%%%%%%%%%%%%%%%%%%%%%%%%%%%%%%%%%%%%%%%%%%%%%%%%%

A perturbation is singular when setting the small parameter equal to zero
changes the mathematical character of the problem.

For example,

\[
\boxed{
\varepsilon y''+y'+y=0
}
\]

is second order when

\[
\varepsilon>0,
\]

but setting

\[
\varepsilon=0
\]

gives

\[
y'+y=0,
\]

which is only first order.

One boundary condition is effectively lost.

This signals a singular perturbation.

%%%%%%%%%%%%%%%%%%%%%%%%%%%%%%%%%%%%%%%%%%%%%%%%%%%%%%%%%%%%
\subsection{Boundary Layers}
\label{subsec:boundary-layers}
%%%%%%%%%%%%%%%%%%%%%%%%%%%%%%%%%%%%%%%%%%%%%%%%%%%%%%%%%%%%

In singular perturbation problems, a regular outer approximation may fail
near a boundary.

A narrow region where the solution changes rapidly is called a boundary
layer.

Suppose its width is

\[
\delta(\varepsilon).
\]

Introduce a stretched variable

\[
\boxed{
X
=
\frac{x-x_0}{\delta(\varepsilon)}.
}
\]

The scale \(\delta(\varepsilon)\) is chosen by dominant balance.

%%%%%%%%%%%%%%%%%%%%%%%%%%%%%%%%%%%%%%%%%%%%%%%%%%%%%%%%%%%%
\subsection{Outer and Inner Solutions}
\label{subsec:outer-inner-solutions}
%%%%%%%%%%%%%%%%%%%%%%%%%%%%%%%%%%%%%%%%%%%%%%%%%%%%%%%%%%%%

The approximation valid away from the boundary layer is called the outer
solution:

\[
\boxed{
u_{\mathrm{out}}.
}
\]

The approximation valid inside the boundary layer is called the inner
solution:

\[
\boxed{
u_{\mathrm{in}}.
}
\]

The inner problem is written using a stretched variable.

Neither approximation alone is usually valid everywhere.

%%%%%%%%%%%%%%%%%%%%%%%%%%%%%%%%%%%%%%%%%%%%%%%%%%%%%%%%%%%%
\subsection{Matched Asymptotic Expansions}
\label{subsec:matched-asymptotic-expansions}
%%%%%%%%%%%%%%%%%%%%%%%%%%%%%%%%%%%%%%%%%%%%%%%%%%%%%%%%%%%%

Matched asymptotic expansions combine inner and outer approximations.

The basic strategy is:

\begin{enumerate}
    \item derive an outer expansion;
    \item identify where it fails;
    \item introduce an inner scaling;
    \item derive an inner expansion;
    \item match the two expansions in an overlap region;
    \item construct a uniformly valid composite approximation.
\end{enumerate}

Symbolically, the matching condition has the form

\[
\boxed{
\text{inner limit}
=
\text{outer limit}.
}
\]

%%%%%%%%%%%%%%%%%%%%%%%%%%%%%%%%%%%%%%%%%%%%%%%%%%%%%%%%%%%%
\subsection{Composite Expansions}
\label{subsec:composite-expansions}
%%%%%%%%%%%%%%%%%%%%%%%%%%%%%%%%%%%%%%%%%%%%%%%%%%%%%%%%%%%%

A common composite approximation is

\[
\boxed{
u_{\mathrm{comp}}
=
u_{\mathrm{out}}
+
u_{\mathrm{in}}
-
u_{\mathrm{overlap}}.
}
\]

The overlap term is subtracted because it is contained in both the inner
and outer approximations.

The result can remain accurate across the entire domain.

%%%%%%%%%%%%%%%%%%%%%%%%%%%%%%%%%%%%%%%%%%%%%%%%%%%%%%%%%%%%
\subsection{Initial Layers}
\label{subsec:initial-layers}
%%%%%%%%%%%%%%%%%%%%%%%%%%%%%%%%%%%%%%%%%%%%%%%%%%%%%%%%%%%%

Boundary-layer behavior can also occur in time-dependent problems.

If a solution changes rapidly near

\[
t=0,
\]

the narrow region is called an initial layer.

A stretched time variable such as

\[
\boxed{
T=\frac{t}{\varepsilon^\alpha}
}
\]

may reveal the fast transient behavior.

%%%%%%%%%%%%%%%%%%%%%%%%%%%%%%%%%%%%%%%%%%%%%%%%%%%%%%%%%%%%
\subsection{Multiple Scales}
\label{subsec:multiple-scales}
%%%%%%%%%%%%%%%%%%%%%%%%%%%%%%%%%%%%%%%%%%%%%%%%%%%%%%%%%%%%

Some perturbation problems contain behavior on several time or spatial
scales.

For example, define

\[
\boxed{
T_0=t,
\qquad
T_1=\varepsilon t,
\qquad
T_2=\varepsilon^2t.
}
\]

Then treat these variables temporarily as independent.

By the chain rule,

\[
\boxed{
\frac{d}{dt}
=
\frac{\partial}{\partial T_0}
+
\varepsilon
\frac{\partial}{\partial T_1}
+
\varepsilon^2
\frac{\partial}{\partial T_2}
+\cdots.
}
\]

This is the basis of the method of multiple scales.

%%%%%%%%%%%%%%%%%%%%%%%%%%%%%%%%%%%%%%%%%%%%%%%%%%%%%%%%%%%%
\subsection{Secular Terms}
\label{subsec:secular-terms}
%%%%%%%%%%%%%%%%%%%%%%%%%%%%%%%%%%%%%%%%%%%%%%%%%%%%%%%%%%%%

A naive perturbation expansion may produce terms such as

\[
\boxed{
\varepsilon t\sin t.
}
\]

Even though \(\varepsilon\) is small, this term becomes order one when

\[
t=O(\varepsilon^{-1}).
\]

Such growing terms are called secular terms.

They cause a regular perturbation expansion to lose validity over long
times.

Multiple-scale methods are designed to remove or control them.

%%%%%%%%%%%%%%%%%%%%%%%%%%%%%%%%%%%%%%%%%%%%%%%%%%%%%%%%%%%%
\subsection{Poincar\'e--Lindstedt Method}
\label{subsec:poincare-lindstedt}
%%%%%%%%%%%%%%%%%%%%%%%%%%%%%%%%%%%%%%%%%%%%%%%%%%%%%%%%%%%%

Weakly nonlinear oscillators often experience a frequency shift.

Instead of expanding only the solution, the Poincar\'e--Lindstedt method
also expands the oscillation frequency.

For example,

\[
\boxed{
\omega
=
\omega_0
+
\varepsilon\omega_1
+
\varepsilon^2\omega_2
+\cdots.
}
\]

A rescaled time variable then absorbs secular growth.

This produces approximations valid over longer time intervals.

%%%%%%%%%%%%%%%%%%%%%%%%%%%%%%%%%%%%%%%%%%%%%%%%%%%%%%%%%%%%
\subsection{WKB Approximation}
\label{subsec:wkb-approximation}
%%%%%%%%%%%%%%%%%%%%%%%%%%%%%%%%%%%%%%%%%%%%%%%%%%%%%%%%%%%%

For differential equations with rapidly varying phases, one often seeks
a solution of the form

\[
\boxed{
y(x)
=
e^{S(x)/\varepsilon}.
}
\]

More generally,

\[
\boxed{
y(x)
\sim
e^{S_0(x)/\varepsilon}
\left(
A_0(x)
+
\varepsilon A_1(x)
+\cdots
\right).
}
\]

This approach is associated with the WKB method.

The letters WKB refer to Wentzel, Kramers, and Brillouin.

%%%%%%%%%%%%%%%%%%%%%%%%%%%%%%%%%%%%%%%%%%%%%%%%%%%%%%%%%%%%
\subsection{WKB for a Second-Order Equation}
\label{subsec:wkb-second-order}
%%%%%%%%%%%%%%%%%%%%%%%%%%%%%%%%%%%%%%%%%%%%%%%%%%%%%%%%%%%%

Consider

\[
\varepsilon^2y''(x)
=
Q(x)y(x).
\]

Assume

\[
y=e^{S/\varepsilon}.
\]

Then

\[
y'
=
\frac{S'}{\varepsilon}e^{S/\varepsilon},
\]

and

\[
y''
=
\left(
\frac{S''}{\varepsilon}
+
\frac{(S')^2}{\varepsilon^2}
\right)
e^{S/\varepsilon}.
\]

Substitution gives

\[
\varepsilon S''
+
(S')^2
=
Q.
\]

The leading equation is

\[
\boxed{
(S_0')^2=Q.
}
\]

Thus

\[
S_0'
=
\pm\sqrt{Q}.
\]

%%%%%%%%%%%%%%%%%%%%%%%%%%%%%%%%%%%%%%%%%%%%%%%%%%%%%%%%%%%%
\subsection{Turning Points}
\label{subsec:turning-points}
%%%%%%%%%%%%%%%%%%%%%%%%%%%%%%%%%%%%%%%%%%%%%%%%%%%%%%%%%%%%

A turning point occurs where

\[
\boxed{
Q(x_0)=0.
}
\]

The standard WKB approximation usually fails near such a point.

A new local scaling must be introduced.

Near a simple turning point, the resulting local equation is often
related to the Airy equation.

This is another example of a singular asymptotic transition.

%%%%%%%%%%%%%%%%%%%%%%%%%%%%%%%%%%%%%%%%%%%%%%%%%%%%%%%%%%%%
\subsection{Special Functions and Asymptotics}
\label{subsec:special-functions-asymptotics}
%%%%%%%%%%%%%%%%%%%%%%%%%%%%%%%%%%%%%%%%%%%%%%%%%%%%%%%%%%%%

Many special functions are defined through

\[
\boxed{
\text{integrals},
\quad
\text{differential equations},
\quad
\text{series}.
}
\]

Exact formulas may be difficult to evaluate for large parameters.

Asymptotic methods provide simpler approximations.

Important examples include

\[
\boxed{
\Gamma(x),
\quad
\text{Bessel functions},
\quad
\text{Airy functions},
\quad
\text{error functions}.
}
\]

%%%%%%%%%%%%%%%%%%%%%%%%%%%%%%%%%%%%%%%%%%%%%%%%%%%%%%%%%%%%
\subsection{Asymptotic Expansions of Integrals by Parts}
\label{subsec:asymptotic-integration-by-parts}
%%%%%%%%%%%%%%%%%%%%%%%%%%%%%%%%%%%%%%%%%%%%%%%%%%%%%%%%%%%%

Repeated integration by parts can generate asymptotic expansions.

Consider

\[
I(x)
=
\int_x^\infty
e^{-t}f(t)\,dt.
\]

Integration by parts gives

\[
I(x)
=
e^{-x}f(x)
+
\int_x^\infty
e^{-t}f'(t)\,dt.
\]

Repeating the process yields a formal expansion involving

\[
f(x),
\quad
f'(x),
\quad
f''(x),
\quad
\ldots
\]

provided the remainder can be controlled.

%%%%%%%%%%%%%%%%%%%%%%%%%%%%%%%%%%%%%%%%%%%%%%%%%%%%%%%%%%%%
\subsection{Optimal Truncation}
\label{subsec:optimal-truncation}
%%%%%%%%%%%%%%%%%%%%%%%%%%%%%%%%%%%%%%%%%%%%%%%%%%%%%%%%%%%%

A divergent asymptotic series can still be extremely accurate when
truncated appropriately.

Often the terms initially decrease in magnitude and later increase.

The best approximation is frequently obtained by truncating near the
smallest term.

This is called optimal truncation.

Thus adding more terms to an asymptotic series can eventually make the
approximation worse.

%%%%%%%%%%%%%%%%%%%%%%%%%%%%%%%%%%%%%%%%%%%%%%%%%%%%%%%%%%%%
\subsection{Exponentially Small Terms}
\label{subsec:exponentially-small-terms}
%%%%%%%%%%%%%%%%%%%%%%%%%%%%%%%%%%%%%%%%%%%%%%%%%%%%%%%%%%%%

Power-series asymptotics may fail to detect terms such as

\[
\boxed{
e^{-1/\varepsilon}
}
\]

as

\[
\varepsilon\to0^+.
\]

Indeed,

\[
\boxed{
e^{-1/\varepsilon}
=
o(\varepsilon^N)
}
\]

for every fixed positive integer \(N\).

Such terms are said to be beyond all algebraic orders.

They become important in refined asymptotic theories.

%%%%%%%%%%%%%%%%%%%%%%%%%%%%%%%%%%%%%%%%%%%%%%%%%%%%%%%%%%%%
\subsection{Stokes Phenomenon}
\label{subsec:stokes-phenomenon}
%%%%%%%%%%%%%%%%%%%%%%%%%%%%%%%%%%%%%%%%%%%%%%%%%%%%%%%%%%%%

In complex asymptotic analysis, exponentially small contributions may
appear or disappear as the argument crosses certain curves in the complex
plane.

This behavior is called the Stokes phenomenon.

It occurs prominently in asymptotic expansions of special functions and
solutions of differential equations.

The phenomenon demonstrates that an asymptotic expansion can depend
subtly on the direction from which a complex limit is approached.

%%%%%%%%%%%%%%%%%%%%%%%%%%%%%%%%%%%%%%%%%%%%%%%%%%%%%%%%%%%%
\subsection{Asymptotics of Sequences}
\label{subsec:asymptotics-of-sequences}
%%%%%%%%%%%%%%%%%%%%%%%%%%%%%%%%%%%%%%%%%%%%%%%%%%%%%%%%%%%%

Asymptotic analysis applies equally to sequences.

For example,

\[
a_n
=
n^2+3n+1
\]

satisfies

\[
\boxed{
a_n\sim n^2.
}
\]

More refined information is

\[
\boxed{
a_n
=
n^2
\left(
1+\frac3n+\frac1{n^2}
\right).
}
\]

Sequence asymptotics are fundamental in combinatorics, probability, and
number theory.

%%%%%%%%%%%%%%%%%%%%%%%%%%%%%%%%%%%%%%%%%%%%%%%%%%%%%%%%%%%%
\subsection{Asymptotics of Binomial Coefficients}
\label{subsec:binomial-asymptotics}
%%%%%%%%%%%%%%%%%%%%%%%%%%%%%%%%%%%%%%%%%%%%%%%%%%%%%%%%%%%%

Using Stirling's formula,

\[
\binom{2n}{n}
=
\frac{(2n)!}{(n!)^2}.
\]

Applying

\[
n!
\sim
\sqrt{2\pi n}
\left(\frac ne\right)^n
\]

gives

\[
\boxed{
\binom{2n}{n}
\sim
\frac{4^n}{\sqrt{\pi n}}.
}
\]

This approximation appears frequently in combinatorics and probability.

%%%%%%%%%%%%%%%%%%%%%%%%%%%%%%%%%%%%%%%%%%%%%%%%%%%%%%%%%%%%
\subsection{Probability and Asymptotics}
\label{subsec:probability-asymptotics}
%%%%%%%%%%%%%%%%%%%%%%%%%%%%%%%%%%%%%%%%%%%%%%%%%%%%%%%%%%%%

Probability theory contains many asymptotic questions.

Examples include

\[
\boxed{
\text{large-sample limits},
}
\]

\[
\boxed{
\text{tail probabilities},
}
\]

\[
\boxed{
\text{rare-event probabilities},
}
\]

and

\[
\boxed{
\text{large deviations}.
}
\]

Laplace-type methods often explain why Gaussian approximations arise in
large systems.

%%%%%%%%%%%%%%%%%%%%%%%%%%%%%%%%%%%%%%%%%%%%%%%%%%%%%%%%%%%%
\subsection{Asymptotic Normality}
\label{subsec:asymptotic-normality}
%%%%%%%%%%%%%%%%%%%%%%%%%%%%%%%%%%%%%%%%%%%%%%%%%%%%%%%%%%%%

A sequence of random variables or estimators may become approximately
normal after suitable centering and scaling.

A typical statement has the form

\[
\boxed{
\frac{X_n-a_n}{b_n}
\Longrightarrow
N(0,1),
}
\]

where \(\Longrightarrow\) denotes convergence in distribution.

This is an asymptotic statement rather than an exact finite-\(n\)
identity.

%%%%%%%%%%%%%%%%%%%%%%%%%%%%%%%%%%%%%%%%%%%%%%%%%%%%%%%%%%%%
\subsection{Asymptotic Analysis in Differential Equations}
\label{subsec:asymptotic-differential-equations}
%%%%%%%%%%%%%%%%%%%%%%%%%%%%%%%%%%%%%%%%%%%%%%%%%%%%%%%%%%%%

Differential equations often contain parameters that make exact
solutions unavailable or uninformative.

Asymptotic analysis can reveal

\[
\boxed{
\text{long-time behavior},
\quad
\text{small-parameter behavior},
\quad
\text{rapid oscillation},
\quad
\text{boundary layers},
\quad
\text{wave propagation}.
}
\]

This makes asymptotic analysis a central tool in applied mathematics.

%%%%%%%%%%%%%%%%%%%%%%%%%%%%%%%%%%%%%%%%%%%%%%%%%%%%%%%%%%%%
\subsection{Uniform and Nonuniform Approximations}
\label{subsec:uniform-nonuniform-asymptotics}
%%%%%%%%%%%%%%%%%%%%%%%%%%%%%%%%%%%%%%%%%%%%%%%%%%%%%%%%%%%%

An approximation may be accurate for each fixed \(x\) but fail to be
uniform over a larger domain.

Suppose

\[
u_\varepsilon(x)
\sim
u_0(x)
\]

for each fixed \(x>0\).

This does not necessarily imply

\[
\boxed{
\sup_x
|u_\varepsilon(x)-u_0(x)|
\to0.
}
\]

Boundary layers are a common source of nonuniformity.

%%%%%%%%%%%%%%%%%%%%%%%%%%%%%%%%%%%%%%%%%%%%%%%%%%%%%%%%%%%%
\subsection{Error Estimates}
\label{subsec:asymptotic-error-estimates}
%%%%%%%%%%%%%%%%%%%%%%%%%%%%%%%%%%%%%%%%%%%%%%%%%%%%%%%%%%%%

An asymptotic formula becomes more informative when accompanied by an
explicit error estimate.

For example,

\[
\boxed{
f(x)
=
g(x)
+
O(x^{-2})
}
\]

provides more information than simply

\[
f(x)\sim g(x).
\]

The error term indicates the scale of the neglected contribution.

Rigorous asymptotic analysis therefore seeks both

\[
\boxed{
\text{leading behavior}
}
\]

and

\[
\boxed{
\text{remainder control}.
}
\]

%%%%%%%%%%%%%%%%%%%%%%%%%%%%%%%%%%%%%%%%%%%%%%%%%%%%%%%%%%%%
\subsection{Common Errors}
\label{subsec:asymptotic-analysis-common-errors}
%%%%%%%%%%%%%%%%%%%%%%%%%%%%%%%%%%%%%%%%%%%%%%%%%%%%%%%%%%%%

\subsubsection{Treating \(\sim\) as Ordinary Equality}

The statement

\[
f\sim g
\]

does not mean

\[
f=g.
\]

It means

\[
\boxed{
\frac fg\to1.
}
\]

%%%%%%%%%%%%%%%%%%%%%%%%%%%%%%%%%%%%%%%%%%%%%%%%%%%%%%%%%%%%
\subsubsection{Confusing Big-O with Asymptotic Equivalence}

From

\[
f=O(g)
\]

one cannot conclude

\[
f\sim g.
\]

For example,

\[
2x=O(x),
\]

but

\[
2x\not\sim x.
\]

%%%%%%%%%%%%%%%%%%%%%%%%%%%%%%%%%%%%%%%%%%%%%%%%%%%%%%%%%%%%
\subsubsection{Confusing Big-O and Little-o}

The statement

\[
f=o(g)
\]

is stronger than

\[
f=O(g).
\]

%%%%%%%%%%%%%%%%%%%%%%%%%%%%%%%%%%%%%%%%%%%%%%%%%%%%%%%%%%%%
\subsubsection{Forgetting the Limiting Regime}

A statement such as

\[
f=O(g)
\]

is incomplete unless the relevant limit is understood.

For example,

\[
x^2=O(x)
\]

is false as \(x\to\infty\), but true as \(x\to0\).

%%%%%%%%%%%%%%%%%%%%%%%%%%%%%%%%%%%%%%%%%%%%%%%%%%%%%%%%%%%%
\subsubsection{Assuming an Asymptotic Series Must Converge}

Many useful asymptotic expansions diverge as infinite series.

Their meaning is based on finite truncation and remainder size.

%%%%%%%%%%%%%%%%%%%%%%%%%%%%%%%%%%%%%%%%%%%%%%%%%%%%%%%%%%%%
\subsubsection{Adding Too Many Terms}

For a divergent asymptotic expansion, terms eventually grow.

Optimal truncation may require stopping before that happens.

%%%%%%%%%%%%%%%%%%%%%%%%%%%%%%%%%%%%%%%%%%%%%%%%%%%%%%%%%%%%
\subsubsection{Ignoring Multiple Scales}

A regular expansion may work for

\[
t=O(1)
\]

but fail when

\[
t=O(\varepsilon^{-1}).
\]

The relevant scale must be identified.

%%%%%%%%%%%%%%%%%%%%%%%%%%%%%%%%%%%%%%%%%%%%%%%%%%%%%%%%%%%%
\subsubsection{Ignoring Boundary Layers}

Setting

\[
\varepsilon=0
\]

in a singular perturbation problem can remove derivatives and boundary
conditions.

The missing behavior is often concentrated in a boundary layer.

%%%%%%%%%%%%%%%%%%%%%%%%%%%%%%%%%%%%%%%%%%%%%%%%%%%%%%%%%%%%
\subsubsection{Using Laplace's Method at the Wrong Point}

For

\[
e^{\lambda\phi(x)},
\]

the dominant region is associated with maxima of

\[
\operatorname{Re}\phi,
\]

not arbitrary critical points.

%%%%%%%%%%%%%%%%%%%%%%%%%%%%%%%%%%%%%%%%%%%%%%%%%%%%%%%%%%%%
\subsubsection{Ignoring Stationary Points}

Rapid oscillations cancel strongly away from points where

\[
\phi'(x)=0.
\]

Stationary points often dominate oscillatory integrals.

%%%%%%%%%%%%%%%%%%%%%%%%%%%%%%%%%%%%%%%%%%%%%%%%%%%%%%%%%%%%
\subsection{Applications}
\label{subsec:asymptotic-analysis-applications}
%%%%%%%%%%%%%%%%%%%%%%%%%%%%%%%%%%%%%%%%%%%%%%%%%%%%%%%%%%%%

\begin{applicationbox}
\textbf{Differential Equations.}
Perturbation methods approximate solutions when small or large parameters
create multiple scales.

\medskip

\textbf{Fluid Mechanics.}
Boundary-layer theory describes rapid changes near solid boundaries.

\medskip

\textbf{Wave Propagation.}
WKB and stationary-phase methods describe rapidly oscillatory waves.

\medskip

\textbf{Quantum Mechanics.}
Semiclassical approximations use small-parameter asymptotics.

\medskip

\textbf{Probability and Statistics.}
Large-sample approximations, tail probabilities, and limiting
distributions rely heavily on asymptotic methods.

\medskip

\textbf{Combinatorics.}
Stirling's formula and saddle-point methods approximate rapidly growing
sequences and coefficients.

\medskip

\textbf{Special Functions.}
Bessel, Airy, Gamma, and related functions have important asymptotic
expansions.

\medskip

\textbf{Fourier and Harmonic Analysis.}
Stationary phase and oscillatory-integral estimates determine
high-frequency behavior.

\medskip

\textbf{Complex Analysis.}
Contour deformation and steepest descent extract dominant contributions
from complex integrals.

\medskip

\textbf{Numerical Analysis.}
Asymptotic estimates guide algorithm design, error prediction, and
parameter selection.

\medskip

\textbf{Mathematical Physics.}
Asymptotic approximations describe systems for which exact formulas are
either unavailable or less informative than limiting behavior.
\end{applicationbox}

%%%%%%%%%%%%%%%%%%%%%%%%%%%%%%%%%%%%%%%%%%%%%%%%%%%%%%%%%%%%
\subsection{Exercises}
\label{subsec:asymptotic-analysis-exercises}
%%%%%%%%%%%%%%%%%%%%%%%%%%%%%%%%%%%%%%%%%%%%%%%%%%%%%%%%%%%%

\begin{exercise}[1]
Show that

\[
x^2+3x+1\sim x^2
\]

as \(x\to\infty\).
\end{exercise}

\begin{exercise}[1]
Show that

\[
\sin x\sim x
\]

as \(x\to0\).
\end{exercise}

\begin{exercise}[1]
Show that

\[
1-\cos x
\sim
\frac{x^2}{2}
\]

as \(x\to0\).
\end{exercise}

\begin{exercise}[1]
Determine whether

\[
x=O(x^2)
\]

as \(x\to\infty\).
\end{exercise}

\begin{exercise}[1]
Determine whether

\[
x^2=o(x)
\]

as \(x\to0\).
\end{exercise}

\begin{exercise}[1]
Explain the difference between

\[
f=O(g)
\]

and

\[
f=o(g).
\]
\end{exercise}

\begin{exercise}[2]
Show that

\[
\log x=o(x^\alpha)
\]

as \(x\to\infty\) for every \(\alpha>0\).
\end{exercise}

\begin{exercise}[2]
Show that

\[
x^m=o(e^x)
\]

as \(x\to\infty\) for every fixed positive integer \(m\).
\end{exercise}

\begin{exercise}[2]
Derive the expansion

\[
\frac{1}{1+x}
=
1-x+x^2+O(x^3)
\]

as \(x\to0\).
\end{exercise}

\begin{exercise}[2]
Find an expansion of

\[
\log(1+x)
\]

through terms of order \(x^3\).
\end{exercise}

\begin{exercise}[2]
Find an expansion of

\[
e^x\cos x
\]

through terms of order \(x^3\).
\end{exercise}

\begin{exercise}[2]
Find the first two terms of the asymptotic expansion of

\[
\sqrt{x^2+x}
\]

as \(x\to\infty\).
\end{exercise}

\begin{exercise}[2]
Use dominant balance to estimate the positive solution of

\[
x^5+x=N
\]

for large \(N\).
\end{exercise}

\begin{exercise}[2]
Use Stirling's formula to approximate \(100!\).
\end{exercise}

\begin{exercise}[2]
Derive

\[
\binom{2n}{n}
\sim
\frac{4^n}{\sqrt{\pi n}}.
\]
\end{exercise}

\begin{exercise}[3]
Apply Laplace's method to

\[
\int_{-1}^{1}
e^{-\lambda x^2}(1+x^2)\,dx
\]

as \(\lambda\to\infty\).
\end{exercise}

\begin{exercise}[3]
Use Watson's lemma to derive an asymptotic expansion of a simple
Laplace transform.
\end{exercise}

\begin{exercise}[3]
Use repeated integration by parts to obtain an asymptotic expansion for

\[
\int_x^\infty
\frac{e^{-t}}{t}\,dt
\]

as \(x\to\infty\).
\end{exercise}

\begin{exercise}[3]
Derive the leading stationary-phase approximation for an oscillatory
integral with one nondegenerate stationary point.
\end{exercise}

\begin{exercise}[3]
Explain why oscillatory integrals decay rapidly when the phase has no
stationary points.
\end{exercise}

\begin{exercise}[3]
Use a regular perturbation expansion to approximate the solution of

\[
x+\varepsilon x^3=1
\]

through first order in \(\varepsilon\).
\end{exercise}

\begin{exercise}[3]
Explain why

\[
\varepsilon y''+y'=0
\]

with two boundary conditions is a singular perturbation problem.
\end{exercise}

\begin{exercise}[3]
Determine an appropriate boundary-layer scaling for a simple
second-order singular perturbation equation.
\end{exercise}

\begin{exercise}[3]
Construct inner and outer approximations for a boundary-layer problem and
state the matching condition.
\end{exercise}

\begin{exercise}[3]
Explain why a secular term such as

\[
\varepsilon t\sin t
\]

causes a perturbation expansion to fail for large \(t\).
\end{exercise}

\begin{exercise}[3]
Introduce multiple time scales

\[
T_0=t,
\qquad
T_1=\varepsilon t
\]

and derive the corresponding first and second time derivatives.
\end{exercise}

\begin{exercise}[3]
Apply the WKB substitution

\[
y=e^{S/\varepsilon}
\]

to

\[
\varepsilon^2y''=Q(x)y
\]

and derive the leading equation for \(S\).
\end{exercise}

\begin{exercise}[4]
Derive Stirling's approximation from the Gamma integral using Laplace's
method.
\end{exercise}

\begin{exercise}[4]
Develop the next correction term in Laplace's method for an integral with
a nondegenerate interior maximum.
\end{exercise}

\begin{exercise}[4]
Study steepest-descent contours for a complex exponential integral.
\end{exercise}

\begin{exercise}[4]
Derive a stationary-phase expansion beyond the leading term.
\end{exercise}

\begin{exercise}[4]
Construct a uniformly valid composite expansion for a singular
perturbation boundary-value problem.
\end{exercise}

\begin{exercise}[4]
Apply the Poincar\'e--Lindstedt method to a weakly nonlinear oscillator.
\end{exercise}

\begin{exercise}[4]
Study the breakdown of WKB approximation near a turning point and derive
the corresponding Airy scaling.
\end{exercise}

\begin{exercise}[4]
Investigate optimal truncation for a divergent asymptotic expansion.
\end{exercise}

\begin{exercise}[4]
Study an example of the Stokes phenomenon for a classical special
function.
\end{exercise}

%%%%%%%%%%%%%%%%%%%%%%%%%%%%%%%%%%%%%%%%%%%%%%%%%%%%%%%%%%%%
\subsection{Conceptual Summary}
\label{subsec:asymptotic-analysis-summary}
%%%%%%%%%%%%%%%%%%%%%%%%%%%%%%%%%%%%%%%%%%%%%%%%%%%%%%%%%%%%

Asymptotic analysis describes mathematical behavior in limiting regimes.

The relation

\[
\boxed{
f\sim g
}
\]

means

\[
\boxed{
\frac fg\to1.
}
\]

Big-O notation gives a magnitude bound,

\[
\boxed{
f=O(g),
}
\]

while little-o expresses asymptotic negligibility,

\[
\boxed{
f=o(g).
}
\]

An asymptotic expansion

\[
\boxed{
f
\sim
a_0\phi_0
+
a_1\phi_1
+
a_2\phi_2
+\cdots
}
\]

does not need to converge as an infinite series.

Its meaning is determined by the error after finite truncation.

Laplace's method analyzes exponentially concentrated integrals.

Stationary phase and steepest descent analyze rapidly oscillatory and
complex integrals.

Stirling's formula gives the fundamental factorial approximation

\[
\boxed{
n!
\sim
\sqrt{2\pi n}
\left(\frac ne\right)^n.
}
\]

Perturbation theory introduces expansions in small parameters.

Regular perturbations preserve the structure of the limiting problem,
while singular perturbations may create

\[
\boxed{
\text{boundary layers},
\quad
\text{multiple scales},
\quad
\text{rapid transitions}.
}
\]

Matched asymptotic expansions combine local approximations from different
scales.

Thus asymptotic analysis provides systematic methods for extracting the
dominant mathematical structure of problems whose exact solutions are
unavailable, complicated, or less useful than their limiting behavior.

%%%%%%%%%%%%%%%%%%%%%%%%%%%%%%%%%%%%%%%%%%%%%%%%%%%%%%%%%%%%
\subsection{Section Connections}
\label{subsec:asymptotic-analysis-connections}
%%%%%%%%%%%%%%%%%%%%%%%%%%%%%%%%%%%%%%%%%%%%%%%%%%%%%%%%%%%%

\chapterconnection{
Asymptotic Analysis builds on Real Analysis through limits, error
estimates, and convergence. Complex Analysis supplies contour deformation,
saddle points, and steepest-descent methods. Fourier and Harmonic Analysis
supply oscillatory integrals and stationary-phase techniques.
Approximation Theory studies approximation within prescribed function
families, while Asymptotic Analysis studies approximations determined by
limiting regimes and scale separation. Differential Equations use regular
perturbation, singular perturbation, boundary layers, multiple scales,
and WKB methods. Probability and Statistics use asymptotic distributions,
large-sample approximations, and Laplace-type estimates. Numerical
Analysis uses asymptotic error formulas to understand computational
methods. The next section, Geometric Analysis, combines analytic methods
with differential-geometric structures and geometric differential
equations.
}

%%%%%%%%%%%%%%%%%%%%%%%%%%%%%%%%%%%%%%%%%%%%%%%%%%%%%%%%%%%%
% SECTION SOURCE TRACKING
%%%%%%%%%%%%%%%%%%%%%%%%%%%%%%%%%%%%%%%%%%%%%%%%%%%%%%%%%%%%

%------------------------------------------------------------
% 5.17 Asymptotic Analysis
%------------------------------------------------------------
%
% Source basis:
%   - Standard asymptotic analysis.
%   - Standard perturbation theory.
%   - Standard asymptotic evaluation of integrals.
%   - Standard singular perturbation methods.
%
% Used for:
%   - asymptotic equivalence;
%   - Big-O and little-o notation;
%   - growth hierarchies;
%   - asymptotic sequences;
%   - asymptotic expansions;
%   - Taylor-based asymptotics;
%   - asymptotic inversion;
%   - dominant balance and scaling;
%   - distinguished limits;
%   - Laplace's method;
%   - endpoint asymptotics;
%   - Watson's lemma;
%   - Gamma function asymptotics;
%   - Stirling's formula;
%   - oscillatory integrals;
%   - nonstationary phase;
%   - stationary phase;
%   - steepest descent;
%   - saddle-point methods;
%   - regular perturbation;
%   - singular perturbation;
%   - boundary and initial layers;
%   - matched asymptotic expansions;
%   - composite expansions;
%   - multiple scales;
%   - secular terms;
%   - Poincare--Lindstedt methods;
%   - WKB approximation;
%   - turning points;
%   - special-function asymptotics;
%   - optimal truncation;
%   - exponentially small terms;
%   - Stokes phenomenon;
%   - sequence and combinatorial asymptotics;
%   - probabilistic asymptotics.
%
% Notes:
%   - Approximation Theory develops best approximation and approximation
%     by prescribed function families.
%   - Complex Analysis develops the contour methods underlying steepest
%     descent and saddle-point analysis.
%   - Harmonic Analysis develops oscillatory-integral estimates in
%     greater analytic depth.
%   - Differential Equations develops perturbation and multiple-scale
%     methods for specific ODE and PDE systems.
%   - Numerical Analysis uses asymptotic error estimates computationally.
%   - Geometric Analysis next combines analysis with geometric structures.
%
%%%%%%%%%%%%%%%%%%%%%%%%%%%%%%%%%%%%%%%%%%%%%%%%%%%%%%%%%%%%